\paragraph{First-order equivariant class on left-invariant backgrounds with torsion: the linearised frame and connection equations (conditional; not closure).}
Take the class with base a four-dimensional Lie group \(G\), vanishing background connection and \(\mathbb D_a=u_a{}^bE_b\), where the \(E_a\) are left-invariant fields with \([E_a,E_b]=\mathsf c_{ab}{}^dE_d\), as in the part on commutant labels at nonzero torsion of the paragraph on one loop about commuting matrices. In section values \(\mathsf R=0\), \(\mathsf F=0\) and \(\mathsf T=-\mathsf c\). This is the torsion of the vanishing connection of the class, not of the Levi-Civita connection of the frame; the convention of the compact toy above, in which the equivariant lifts of a torsion-free connection have \(\mathsf T=0\), is unchanged. The background solves the matrix equation when \(\eta^{ac}\mathsf c_{ab}{}^d\mathsf c_{cd}{}^e=0\). Let
\[
\mathsf F_{ab}[A]=E_aA_b-E_bA_a-\mathsf c_{ab}{}^dA_d,\qquad
\mathcal Q_b[A]=E^a\mathsf F_{ab}[A]+\mathsf c_{ab}{}^c\mathsf F^a{}_c[A],
\]
the linearised third bracket, that is the equation of the order-zero field \(A_a\) of the class on this background, which is invariant under \(A_b\to A_b+E_b\Lambda\). Let \(\rho_I\) be the right-invariant fields on \(G\), which commute with all \(\mathbb D_a\) and \(\mathcal S_{cd}\), and write \(E_c=\mathsf L_c{}^I\rho_I\), with \(\mathsf L\) an invertible matrix of functions. For \(\mathbb D_a\to\mathbb D_a-h_a{}^c\mathbb D_c+\tfrac12\omega_a{}^{cd}\mathcal S_{cd}\) with arbitrary functions \(h\) and \(\omega\) on \(G\), the linearised first two brackets of the matrix equation are
\[
\mathcal E^{(2)}{}_b{}^{de}=\mathcal Q_b[\omega^{de}],\qquad
\mathcal E^{(1)}{}_b{}^e\big|_{\omega=0}=\mathcal Q_b[f^I]\,(\mathsf L^{-1})_I{}^e,\qquad
f_b{}^I=-h_b{}^c\mathsf L_c{}^I,
\]
with \(\omega^{de}\) the one-form with components \(\omega_a{}^{de}\). The first holds because, about this background, the linearised curvature is \(\mathsf R_{ab}{}^{de}=\mathsf F_{ab}[\omega^{de}]\) and contains no \(h\); the second is the adjoint compatibility of the paragraph on the pairing of the action, for \(\mathcal C=\rho_I\). Both were checked symbolically for arbitrary \(h\) and \(\omega\) on flat space, on the null-rotation algebra and the second algebra of the paragraph on one loop about commuting matrices, and on the Heisenberg algebra below. So on every such background the forty equations are block-triangular, the connection sourcing the frame, and each of the six components \(\omega^{de}\) and each of the four frame labels \(f^I\) obeys the equation \(\mathcal Q\) of the order-zero field. In particular the Lorentz modes that stop being commutant at \(\mathsf c\neq0\) obey the gauge-invariant equation \(\mathcal Q\); no term breaking \(A_b\to A_b+E_b\Lambda\) appears. At generic characters the kernel of the forty equations is the ten gauge directions (numerical, on the three algebras with \(\mathsf c\neq0\)). Whether such a background gaps the torsion modes is therefore the question whether it gaps \(\mathcal Q\). On flat space and on the null-rotation background in the gauge with constant frame below, moreover, every constant matrix \(G\) for which \(G\) times the forty linearised equations is formally self-adjoint vanishes on the frame equations (numerical), so no constant recombination of these equations is the Euler--Lagrange system of a quadratic Lagrangian in \(h\) and \(\omega\) with constant coefficients; the ten Maxwell equations of the paragraph on the pairing of the action are in variables that contain \(y\) explicitly.

The null-rotation algebra solves the equation for \(\eta\) and not for the Euclidean contraction. The gauge with constant frame and \(\omega_a{}^{12}=\vartheta k_a\), \(k_a=(1,0,0,1)\), has the same torsion for \(\vartheta=-\tfrac12\), with \(\mathsf T_{abc}\mathsf T^{abc}=0\). There the characteristic roots of the forty equations in the character sectors lie on the real cones \((q+j\vartheta k)^2=0\), with multiplicities \(4\), \(6\) and \(2\) for \(j=0\), \(j=\pm1\) and \(j=\pm2\), twenty in all as at the flat point, and with chains of length at most three (roots numerical at \(\vartheta=-\tfrac12,\tfrac12,1,2\); chains numerical and ranks exact in rational arithmetic at sampled points for \(\vartheta=1\)). So these sectors have neither a gap nor a complex root. The Heisenberg algebra \([E_3,E_1]=\kappa E_2\), with \(\kappa>0\) and with \(E_0\) and \(E_2\) central, solves the equation for both contractions, since \(\mathrm{ad}_a\mathrm{ad}_c=0\) for all \(a,c\). Its frame metric is that of time times the Heisenberg group, and its torsion has \(\mathsf T_{abc}\mathsf T^{abc}=2\kappa^2\), vanishing trace and a nonzero totally antisymmetric part, the part read as the field strength of a \(B\)-field in~\cite{FurutaHanadaKawaiKimura2007}. In coordinates with \(E_0=\partial_0\), \(E_1=\partial_1+\tfrac12\kappa y^3\partial_2\), \(E_2=\partial_2\) and \(E_3=\partial_3-\tfrac12\kappa y^1\partial_2\), let \(p>0\), \(\phi=\exp\bigl(\mathrm iq_0y^0+\mathrm ipy^2-\tfrac14\kappa p((y^1)^2+(y^3)^2)\bigr)\) and \(A=(0,\phi,0,\mathrm i\phi)\). Then, exactly and for \(\eta\),
\[
\mathcal Q_b[A]=(q_0^2-p^2+\kappa p)\,A_b,\qquad\mathsf F_{32}[A]=p\,\phi,
\]
so \(A\), which is not a gauge direction, solves \(\mathcal Q[A]=0\) at \(q_0^2=p(p-\kappa)\), which is negative for \(0<p<\kappa\). At \(q_0=-\mathrm i\Gamma\), \(\Gamma=\sqrt{p(\kappa-p)}\le\kappa/2\) with equality at \(p=\kappa/2\), the real part of \(A\) is a real solution growing like \(\mathrm e^{\Gamma y^0}\); it is Gaussian in \(y^1,y^3\), and its superpositions over \(p\) in a closed subinterval of \((0,\kappa)\) are square-integrable on each slice of constant \(y^0\). By the reduction, for each right-invariant label \(J\) the frame fluctuation with \(f_b{}^I=A_b\delta^I{}_J\) solves the linearised first two brackets with nonzero linearised torsion (exact), and in the sector of fixed \(p\) the lowest level of the forty equations has ten times the multiplicity of that of \(\mathcal Q\) (numerical). In that sector the roots of \(\mathcal Q\) are \(q_0^2=p^2+\kappa p(2n+1\mp2)\), \(n\ge0\) (numerical), which has the form of the Landau levels of a charged vector field in a homogeneous magnetic field, with \(\kappa p\) in the place of the field; for that field an unstable mode is found in~\cite{NielsenOlesen1978}. For the Euclidean contraction the same \(A\) gives \(\mathcal Q[A]=-(q_0^2+p^2-\kappa p)A\), which vanishes at real \(q_0\) for \(0<p<\kappa\); no form is used, so no sign is attached to this.

These statements are for Equation~\eqref{eq:Smatrix} with \(C_{\mu\nu\lambda}=0\) and no matter term, tree level, linear order in the fluctuation, left-invariant backgrounds with vanishing connection on a four-dimensional group and the contractions named, with the spectra stated only for the two unimodular examples. The growing modes are solutions of the linear equations; their nonlinear evolution, the regulated trace, the domains of the unbounded operators, loops and a norm are not treated. Nothing here shows that the matrix integral excludes the Heisenberg background, selects this class or its algebra \(\mathcal A\), or yields the frame and its non-degeneracy, so Assumptions~\ref{ass:frame-lock}, \ref{ass:nondeg}, \ref{ass:lorentzian}, \ref{ass:local-lorentz} and~\ref{ass:EFT-a2} are not discharged; Problems~D, F and~N remain open and PROGRAMME-CONDITIONAL.
